\documentclass[10pt,a4paper]{amsart}
\usepackage[margin=19mm]{geometry}
\usepackage{amsmath,amssymb,mathtools}
\usepackage[scr=rsfs,cal=euler]{mathalfa}
\usepackage{xfrac}
\usepackage[unicode,hidelinks,bookmarks=false]{hyperref}
\newcommand{\ie}{i.e.\ }
\newcommand{\cf}{cf.\ }
\newcommand{\resp}{respectively\ }
\newcommand{\Subsection}{\subsection}
\providecommand{\nobreakdash}{\nobreak\hbox{-}}
\setlength{\emergencystretch}{2em}
\multlinegap0pt
\usepackage{amssymb}
\usepackage{xcolor}
\usepackage{algorithm}\usepackage{algpseudocode}
\usepackage[shortlabels]{enumitem}
\newtheorem{lemma}{Lemma}
\newcommand{\supp}{\operatorname{supp}}
\title{On the continuity of the product of distributions in local Sobolev spaces}
\author{Stefan Tuti\'c}
\date{Source: arXiv:2512.23521v1 (CC BY 4.0)}
\begin{document}
\maketitle

\noindent\emph{Source and licence.} On the continuity of the product of distributions in local Sobolev spaces. Version arXiv:2512.23521v1 (CC BY 4.0), distributed by arXiv under the Creative Commons Attribution 4.0 International licence, http://creativecommons.org/licenses/by/4.0/, as recorded in the arXiv licence metadata for this exact version and shown on the abstract page https://arxiv.org/abs/2512.23521v1. Full licence notice: \texttt{licenses/arxiv-2512.23521-CC-BY-4.0.txt}.\par\smallskip

\noindent\textit{Editorial scope.} Counted unit: the article's lemma WF\_tenz\_proiz (source0.tex lines 122--133). The article's complete original proof of it is delivered with it (source0.tex lines 134--214, one \texttt{proof} environment of 81 lines). No mathematics is written, repaired or re-derived: the statement and the proof are the article's own lines, in the article's own notation, and no retained line is substituted or re-flowed.  No further source-local context is needed: every cross-reference of every retained line resolves inside the two delivered spans.  Nothing was omitted from any retained span, so the package carries no omission record.  No retained line contains a citation, so the wrapper carries no bibliography and no bibliography entry.  No retained line contains a figure environment, an image-inclusion command or drawing code, so no image is delivered and the package declares no asset dependency.  Editorial formatting only, in addition: an AMS \texttt{amsart} preamble that restates the article's own declarations and macros used by the retained lines, namely the environments \texttt{lemma} on separate counters, together with the standard packages those lines need.  The retained environments print in this document as Lemma 1 in order of appearance, whereas the article prints the same items with the numbering of its own full document.  The complete native source of the article as distributed in the arXiv e-print for this exact version is bundled unchanged as \texttt{sources/191809/source0.tex}.
\begin{lemma}
\label{WF_tenz_proiz}
Let $O\subseteq \mathbb{R}^m$ and $U\subseteq \mathbb{R}^n$ be open sets. Let $r,r_1,r_2,r',r''\in\mathbb{R}$ satisfy $r \leq r_1 + \min \{0,r''\}$ and $r \leq r_2 + \min \{0,r'\}$. For all $u\in H_{loc} ^{r'} (O), v\in H_{loc} ^{r''} (U)$ it holds that
\begin{align}
\label{talsnifrontinkl}
\begin{split}
WF^r (u\otimes v)
&\subseteq (WF^{r_1} (u) \times WF^{r_2} (v)) \cup (WF^{r_1} (u) \times (\supp v \times \{0\})) \\
&\quad \cup ((\supp u \times \{0 \}) \times WF^{r_2} (v)).
\end{split}
\end{align}
\end{lemma}

\begin{proof}
We start off by showing
\begin{align}
\label{sigma}
\Sigma^r (u\otimes v)
\subseteq
(\Sigma^{r_1} (u) \times \Sigma^{r_2} (v)) \cup (\Sigma^{r_1} (u) \times \{0\}) \cup (\{0\} \times \Sigma^{r_2} (v)),
\end{align}
for $u\in \mathcal{E}' (O) \cap H_{loc} ^{r'} (O) , v\in \mathcal{E}' (U) \cap H_{loc} ^{r''} (U)$. Let $(\xi_0,\eta_0)\notin (\Sigma^{r_1} (u) \times \Sigma^{r_2} (v)) \cup (\Sigma^{r_1} (u) \times \{0\}) \cup (\{0\} \times \Sigma^{r_2} (v))$ satisfies $(\xi_0,\eta_0)\neq (0,0)$. We distinguish three cases:
\begin{enumerate}
\item \underline{$\xi_0 = 0$:}

From $(\xi_0,\eta_0)\notin \{0\} \times \Sigma^{r_2} (v)$, that is $\eta_0 \notin \Sigma^{r_2} (v)$, we get there exists an open cone $V$ which is a neighborhood of $\eta_0$ such that
\begin{align}
\label{eta0}
\lVert \langle\cdot \rangle^{r_2} \mathcal{F} v \rVert_{L^2 (V)}
< \infty.
\end{align}
Set $W:=\{(\xi , \eta) \in \mathbb{R}^{m+n} \mid \eta\in V , |\xi| < |\eta| \}$. Clearly, $(\xi_0,\eta_0)\in W$. If $r'\geq 0$ then $\hat{u}\in L^2 (\mathbb{R}^m)$, because $u\in H_{loc}  ^{r'} (O)$ and has a compact support. Therefore,
\begin{align*}
\lVert \langle \cdot \rangle^{r} \mathcal{F} (u\otimes v)  \rVert_{L^2 (W)} ^2
&= \iint_{\substack{ \eta\in V \\ |\xi|<|\eta| }} \langle (\xi ,\eta) \rangle^{2r} |\mathcal{F}u (\xi)|^2 |\mathcal{F}v (\eta)|^2 \mathrm{d}\xi \mathrm{d}\eta \\
\boxed{\substack{\langle (\xi ,\eta) \rangle^{2r} \leq c \langle \eta \rangle^{2r}, \\ \text{for } |\xi| < |\eta| }}
&\leq c \iint_{\substack{ \eta\in V \\ |\xi|<|\eta| }} \langle \eta \rangle^{2r} |\mathcal{F}u (\xi)|^2 |\mathcal{F}v (\eta)|^2 \mathrm{d}\xi \mathrm{d}\eta\\
\boxed{\substack{ r \leq r_2 + \min \{0,r'\} = r_2 \\ }}
&\leq c \iint_{\substack{ \eta\in V \\ |\xi|<|\eta| }} \langle \eta \rangle^{2r_2} |\mathcal{F}u (\xi)|^2 |\mathcal{F}v (\eta)|^2 \mathrm{d}\xi \mathrm{d}\eta\\
&\leq c \iint_{\substack{ \eta\in V \\ \xi \in \mathbb{R}^m }} \langle \eta \rangle^{2r_2} |\mathcal{F}u (\xi)|^2 |\mathcal{F}v (\eta)|^2 \mathrm{d}\xi \mathrm{d}\eta\\
&= c\int_{\mathbb{R}^m} |\mathcal{F}u (\xi)|^2 \mathrm{d}\xi \int_{\eta\in V} \langle \eta \rangle^{2r_2} |\mathcal{F}v (\eta)|^2 \mathrm{d}\eta
<\infty.
\end{align*}
Now, let $r'<0$. When $r_2 \geq 0 $, and $|\xi|\leq |\eta|$, it follows
\begin{align}
\label{triv}
\langle ( \xi , \eta ) \rangle^r
\leq \langle ( \xi , \eta ) \rangle^{r_2} \langle ( \xi , \eta ) \rangle^{r'}
\leq c \langle \eta \rangle^{r_2} \langle \xi \rangle^{r'}.
\end{align}
For $r_2 < 0$, \eqref{triv} trivially holds. Therefore,
\begin{align*}
\lVert \langle \cdot \rangle^{r} \mathcal{F} (u\otimes v) \rVert_{L^2 (W)} ^2
&\leq c^2 \iint_{\substack{ \eta\in V \\ |\xi|<|\eta| }} \langle \eta \rangle^{2r_2} \langle \xi \rangle^{2r'} |\mathcal{F}u (\xi)|^2 |\mathcal{F}v (\eta)|^2 \mathrm{d}\xi \mathrm{d}\eta\\
&\leq c^2\int_{\xi\in\mathbb{R}^m} \langle \xi \rangle^{2r'} |\mathcal{F}u (\xi)|^2 \mathrm{d}\xi \int_{\eta\in V } \langle \eta \rangle^{2r_2} |\mathcal{F}v (\eta)|^2 \mathrm{d}\eta
< \infty,
\end{align*}
since $u\in H_{loc}  ^{r'} (O)$ and has a compact support.

\item \underline{$\eta_0 = 0$:}

The proof is analogous to the previous case.

\item \underline{$\xi_0\neq0 , \eta_0\neq 0 $:}

Since $(\xi_0, \eta_0)\notin \Sigma^{r_1} (u) \times \Sigma^{r_2} (v)$, lets assume that $\eta_0 \notin \Sigma^{r_2} (v)$. This means that \eqref{eta0} holds true. Choose $R>\frac{|\xi_0|}{|\eta_0|}$ and notice that $(\xi_0 , \eta_0) \in \{ (\xi , \eta)\in \mathbb{R}^{m+n} \mid \eta\in V , |\xi| < R |\eta | \} =: W$. Now, arguing in the same manner as in the first case, we get that
\begin{align*}
\lVert \langle \cdot\rangle^{r} \mathcal{F} (u\otimes v) \rVert_{L^2 (W)} < \infty.
\end{align*}
The case $\xi_0 \notin \Sigma^{r_1} (u)$ follows analogously.
\end{enumerate}
For $u\in \mathcal{D}' (O), v\in \mathcal{D}' (U), (x_0 , y_0) \in O\times U$ and $s\in\mathbb{R}$ it holds that
\begin{align}
\label{sigmaxy}
\begin{split}
&\Sigma_{(x_0 , y_0)} ^s (u\otimes v)\\
&= \bigcap \{ \Sigma^s ( (\varphi u) \otimes (\psi v) ) \mid \varphi \in \mathcal{D} (O) , \varphi (x_0) \neq 0 , \psi\in\mathcal{D} (U) , \psi (y_0) \neq 0\}.
\end{split}
\end{align}
This follows from the fact that $\Sigma^s (\rho w) \subseteq \Sigma^s (w)$ for $\rho \in \mathcal{D} (O) , w \in \mathcal{E}' (O)$. To see this, assume that $(\xi_0 , \eta_0)$ belongs to the right hand side of (\ref{sigmaxy}). Then for arbitrary $\phi\in\mathcal{D} (O\times U) , \phi (x_0 , y_0) \neq 0$ we can choose open sets $M\subseteq O , N\subseteq U$ such that $(x_0,y_0)\in M\times N$ and $\phi \neq 0$ on $M\times N$. Pick $\varphi\in \mathcal{D} (M)$ and $\psi\in\mathcal{D}(N)$ such that $\varphi (x_0) \neq 0$ and $\psi (y_0) \neq 0$. As $\varphi \otimes \psi = (\varphi \otimes \psi)\phi/\phi$, we infer $(\xi_0 , \eta_0) \in \Sigma^s ((\varphi u) \otimes (\psi v)) = \Sigma^s (\phi ((\varphi \otimes \psi)/\phi) (u\otimes v)) \subseteq \Sigma^s (\phi (u\otimes v)) $. The other inclusion follows immediately from the definition of the set $\Sigma_{(x_0 , y_0)} ^s (u\otimes v)$.

Now, let $u\in H_{loc} ^{r'} (O), v\in H_{loc} ^{r''} (U)$. From (\ref{sigma}) and (\ref{sigmaxy}) we infer
\begin{align*}
\Sigma_{(x_0 , y_0)} ^r (u\otimes v)
\subseteq (\Sigma_{x_0 } ^{r_1} (u) \times \Sigma_{y_0 } ^{r_2} (v)) \cup (\Sigma_{x_0 } ^{r_1} (u) \times \{0 \}) \cup ( \{0 \} \times \Sigma_{y_0 } ^{r_2} (v)),
\end{align*}
and therefore
\begin{align*}
WF^r (u\otimes v)
&= \bigcup_{\substack{ x\in\supp u \\ y\in \supp v }} \{(x,y)\} \times \Sigma_{(x,y)} ^r (u\otimes v)\\
&\subseteq (WF^{r_1} (u) \times WF^{r_2} (v)) \cup (WF^{r_1} (u) \times (\supp v \times \{0\}))\\
&\quad \cup ((\supp u \times \{0 \}) \times WF^{r_2} (v)).
\end{align*}
\end{proof}

\end{document}
